\pagenumbering{gobble}

\name{刘波}

\contactInfo{杭州}{15 年系统级开发经验}{C++ / Rust / Go}{kumustone@gmail.com}{18143431580}{硕士（中国海洋大学 · 通信与信息系统）}{1985 年生}

\section{个人简介}

\sloppy
拥有基础设施、网关与安全领域深厚开发背景，精通网络编程与 HTTP 网关架构，具备百万级并发及分布式系统设计经验，持有 CISSP 认证并兼具甲乙方安全视角。擅长从技术架构到团队管理的全链路把控，既能主导高可用微服务架构的设计落地，也能带领团队高效交付核心业务。拥有基础设施、网关与安全领域深厚开发背景，精通网络编程与 HTTP 网关架构，具备百万级并发及分布式系统设计经验，持有 CISSP 认证并兼具甲乙方安全视角。擅长从技术架构到团队管理的全链路把控，既能主导高可用微服务架构的设计落地，也能带领团队高效交付核心业务。

\vspace{0.2cm}

\section{核心能力}
\begin{itemize}[parsep=0.2ex]
    \setstretch{1.1}
  \item \textbf{底层系统开发}: 精通 C++/Rust 底层开发，熟悉 Linux 网络与 epoll 等高并发 I/O 模型，具备高性能网络库与基础组件研发能力。
  \item \textbf{网关与实时通信}: 精通 L7 代理与网关架构，拥有 Nginx 模块化开发与 C++ 扩展经验；熟练掌握 TCP 长连接、会话管理与实时推送链路，保障高并发场景下的稳定性与可观测性。
  \item \textbf{性能与高可用保障}: 拥有容量规划、全链路压测与线上问题根因定位经验，擅长解决 CPU 瓶颈、GC、调度、内存泄漏等问题，保障高吞吐与低延迟。
  \item \textbf{安全与业务风控体系}: 具备 WAF 引擎、策略编排与强对抗治理能力，持有 CISSP 证书及甲乙方安全经验；主导风控中台 0-1 建设，构建流量‑特征‑规则‑处置闭环，精通反作弊、反欺诈、反洗钱、IM 风控等业务风险治理。
\end{itemize}

\vspace{0.2cm}

\section{工作经验}
\vspace{0.2cm}

\datedsubsection{闪捷信息科技有限公司}{系统架构 / 高并发网关}{2022.04-至今}

\thirdheading{高性能安全代理与审计网关 (Nginx + C++)}
\begin{itemize}
    \setstretch{1.1}
    \item \textbf{架构设计}: 自研 Nginx 核心扩展模块，基于 C++ 实现请求、响应头、响应体全生命周期流量管控与代理逻辑。
    \item \textbf{安全能力建设}: 实现 HTTP 反向代理全链路访问控制、数据脱敏、账号发现、接口梳理与审计能力，支撑高并发场景稳定运行。
    \item \textbf{性能优化}: 纯代理模式 QPS 15w/s，业务模式 QPS 3w/s；自研高性能 JSON 解析器与账号识别模块，解决高并发下正则导致的 CPU 瓶颈问题。
    \item \textbf{工程化落地}: 通过敏捷开发与 CI 持续集成快速迭代，支撑千万级项目 POC 与交付落地；服务金融、政府等头部客户，落地告警抑制、异常流量回退与全链路压测体系。
\end{itemize}

\thirdheading{核心敏感数据识别引擎 (Rust 核心模块)}
\begin{itemize}
    \setstretch{1.1}
    \item \textbf{架构与功能}: 基于 Rust 自研高性能敏感数据识别与脱敏底座，支持 80+ 种敏感数据类型，覆盖结构化、半结构化及文件数据，以统一 API 能力支撑网关、脱敏服务、Web 展示、离线 SDK 等多场景复用。
    \item \textbf{性能与价值}: 核心模块性能提升 10 倍，整体链路吞吐提升 1 倍以上，实现多业务规模化落地，显著降低接入成本与交付周期。
\end{itemize}

\thirdheading{Nexis 下一代异步安全网关 (Rust) — 核心负责人 / 架构师}
\begin{itemize}
    \setstretch{1.1}
    \item \textbf{架构重构}: 针对 Nginx 进程模型瓶颈，基于 Pingora（Rust async）重构异步网关底座，采用插件化架构承载安全策略，实现全链路能力动态挂载。
    \item \textbf{性能成果}: QPS 提升 138\%，RT 降低 63\%，稳定支撑复杂安全逻辑，可扩展支持 AI 网关实时审计场景。
\end{itemize}

\vspace{0.2cm}

\datedsubsection{玩物得志}{安全专家 / 团队 Leader}{2019.08-2022.03}

\thirdheading{业务风控体系从 0-1 建设 (核心架构师)}
\begin{itemize}
    \setstretch{1.1}
    \item \textbf{对抗体系建设}: 围绕注册、登录、支付、提现等核心链路构建特征工程与规则体系，涵盖设备指纹、行为序列、人机识别、反爬虫、反作弊等能力。
    \item \textbf{风险治理}: 主导 IM 反诈、反洗钱、私域导流等专项治理，通过实时风控与离线审计闭环，有效降低业务资损与合规风险。
\end{itemize}

\thirdheading{高性能文本反垃圾系统 (自研)}
\begin{itemize}
    \setstretch{1.1}
    \item \textbf{策略编排}: 对接多家第三方风控能力，实现可插拔式接入与策略编排，支撑业务快速迭代。
    \item \textbf{自研与优化}: 基于 DFA、正则与自定义规则实现引擎，通过多级缓存与规则预编译提升性能，单机吞吐量提升 3 倍；核心场景替代商业服务，年度安全成本降低 60\%。
\end{itemize}

\thirdheading{企业级安全治理与等保合规}
\begin{itemize}
    \setstretch{1.1}
    \item \textbf{合规建设}: 主导完成等保三级测评，构建关键链路审计体系，收敛风险暴露面，建立常态化漏洞治理机制。
    \item \textbf{安全工程化}: 落地设备指纹、接口防护、客户端加固等基础安全能力，从架构层面降低爬虫、注入等安全风险。
\end{itemize}

\vspace{0.2cm}

\datedsubsection{蘑菇街}{资深开发工程师}{2015.05-2019.08}

\thirdheading{蘑菇街 自研IM}
\begin{itemize}
    \setstretch{1.1}
    \item 自研 C++ 跨平台长连接 SDK，基于 epoll/kqueue 实现异步非阻塞 I/O，提供分包组包、状态机、自动重连、加密与网络诊断能力，成为公司通用长连接基础组件，降低多端接入成本。
    \item 负责 IM2.0 接入服务、登录鉴权、会话管理及 AppServer 模块全流程开发与上线，保障高并发长连接稳定性。
\end{itemize}

\thirdheading{蘑菇街 WAF网关}
\begin{itemize}
    \setstretch{1.1}
    \item \textbf{架构设计}: 自研全站 WAF 网关，实现网关与检测规则解耦，通过模块化规则链完成实时攻防拦截，支持检测服务水平扩展。
    \item \textbf{性能优化}: 深度优化协议与调用链路，将检测耗时控制在毫秒级，满足高并发场景 10ms 超时要求。
    \item \textbf{落地成果}: 支撑全站流量接入，单机峰值 QPS 20K+；爬虫高峰期超时率从 10\% 降至 0.01\%，高并发场景下稳定性显著提升。
\end{itemize}

\thirdheading{企业级信息安全治理}
\begin{itemize}
    \setstretch{1.1}
    \item \textbf{安全平台}: 研发自动化漏洞扫描平台，覆盖主机、中间件、代码静态扫描，建立漏洞闭环治理流程。
    \item \textbf{主机安全}: 研发 HIDS Agent 与服务端，实现命令审计、文件监控、基线检查；基于 ipset/iptables 实现动态 ACL 防火墙，强化主机安全防护。
    \item \textbf{攻防演练}: 参与渗透测试与应急响应，将攻防成果反哺安全平台与防护规则。
\end{itemize}

\vspace{0.2cm}

\datedsubsection{北京中交兴路车联网}{开发经理 / 高级软件工程师}{2012.04-2015.05}

\begin{itemize}
    \setstretch{1.1}
    \item 拥有 1 项国家发明专利（事件驱动模型数据透传服务），主导核心服务集群化与容灾改造，服务可用性大幅提升。
    \item 设计多级缓存与数据一致性方案，显著降低设备上线异常率，提升系统稳定性。
    \item 自研 epoll 异步网络库与一致性哈希 Redis 客户端，通过锁优化实现单进程并发 5w→25w，消息转发能力达 7w/s。
\end{itemize}

\vspace{0.2cm}

\datedsubsection{华为技术有限公司}{开发工程师}{2011.05-2012.04}

\begin{itemize}
    \setstretch{1.1}
    \item 负责虚拟文件管理（VFM）模块研发，保障文件系统在严苛环境下的高稳定性。
    \item 深入 Linux 内核与 C 语言底层，定位并解决内存泄漏、死锁等高级故障，建立高可靠软件设计与工程化交付理念。
    \item 技术栈：Linux、C、FSM、事件驱动、VSFTP。
\end{itemize}

\vspace{0.2cm}

\datedsubsection{北京千方集团科技有限公司}{开发工程师}{2009.10-2011.05}

\begin{itemize}
    \setstretch{1.1}
    \item 基于 Socket、epoll 研发高并发接入模块，支撑海量 TCP 长连接接入，深度解析 GPS 808/809 国标协议。
    \item 技术栈：C++、线程池、LVS、高性能异步 I/O。
\end{itemize}

\vspace{0.2cm}

\section{技术栈 \& 证书}
\begin{itemize}[parsep=0.2ex]
    \setstretch{1.1}
  \item \textbf{核心语言}: \textbf{C/C++}（精通，10 年+经验）, \textbf{Rust}（3 年）, \textbf{Go}（使用相对少，约 3 年）, Python
  \item \textbf{系统底座}: Pingora, Nginx (C Module), Linux Kernel (Netfilter), Docker / K8s
  \item \textbf{数据组件}: Redis, Kafka, ClickHouse, MySQL, Zookeeper
  \item \textbf{专业资质}: \textbf{CISSP}（国际注册信息系统安全专家）, CET-6, 国家发明专利持有人
  \item \textbf{荣誉奖励}: 卫星导航定位科学技术个人一等奖，优秀员工
  \item \textbf{核心标签}: 高性能网关 (L7 Proxy / Nginx Module / 高并发架构)；内核级优化 (应用层零拷贝 / 内存池化 / 协议栈调优)；安全防护引擎 (WAF 自研 / 实时风控 / HIDS)；复杂故障根因 (内存泄漏排查 / GC 抖动治理 / 性能压测)
\end{itemize}

\vspace{0.2cm}

\section{教育背景}
\begin{tabularx}{\textwidth}{@{\extracolsep{\fill}}X X X X@{}}
  \textbf{中国海洋大学} & \textit{通信与信息系统} & \textit{硕士} & \textit{2007.9 - 2010.6} \\
  \textbf{怀化学院} & \textit{电子信息科学与技术} & \textit{本科} & \textit{2003.9 - 2007.6} \\
\end{tabularx}

